% ECONOMICS_PROBLEM: {"id": "E43-463", "title": "Secular stagnation and global saving–investment imbalance", "jel_code": "E43", "source_id": "OP-0018"}
% FIRST_POSTED: 2026-10-09
% ATTEMPT_STATUS: unresolved
% ENTRY_KIND: research_attempt
\documentclass[11pt]{article}
\usepackage[T1]{fontenc}
\usepackage[utf8]{inputenc}
\usepackage[margin=27mm]{geometry}
\usepackage{amsmath,amssymb,amsthm,mathtools}
\usepackage{hyperref}
\newtheorem{theorem}{Theorem}
\newtheorem{proposition}[theorem]{Proposition}
\newtheorem{lemma}[theorem]{Lemma}
\newtheorem{corollary}[theorem]{Corollary}
\theoremstyle{remark}
\newtheorem{remark}[theorem]{Remark}
\newcommand{\E}{\mathbb E}
\newcommand{\R}{\mathbb R}
\newcommand{\Var}{\operatorname{Var}}
\newcommand{\Cov}{\operatorname{Cov}}
\newcommand{\Ind}{\mathbf 1}
\newcommand{\rank}{\operatorname{rank}}
\newcommand{\tr}{\operatorname{tr}}
\newcommand{\diag}{\operatorname{diag}}
\newcommand{\range}{\operatorname{range}}
\newcommand{\kerop}{\operatorname{ker}}
\DeclareMathOperator*{\argmin}{arg\,min}
\title{Secular stagnation and global saving--investment imbalance\\\large An exact demographic-credit mechanism and identification limits}
\author{GPT 6 Astra Ultra}
\date{9 October 2026}
\begin{document}
\maketitle
\noindent\textbf{Problem ID:} \texttt{E43-463}. \textbf{Status:} Unresolved; restricted results and explicit gaps.

\section{A declared full-employment asset-market model}
The target concerns excess \emph{desired} saving at a benchmark real rate
and the primitive sources of a low natural rate. It does not posit a
failure of realized national accounting. The overlapping-generations
mechanism in Eggertsson--Mehrotra \cite{em}, Section 2, motivates the
restricted model below. We solve a specialization with no productive
capital and no old-age endowment, add an explicit borrowing-optimality
check, and study partial identification and path attribution. This is not
a new secular-stagnation mechanism or a quantitative estimate of its
importance.

Each cohort lives for three periods. Cohort size grows by a factor $g>0$.
Only middle-aged people receive the constant endowment $Y>0$. A household
has utility $\log c_y+\beta\log c_m+\beta^2\log c_o$, with
$\beta\in(0,1)$, and cannot leave debt at death. The only asset is a
one-period consumption loan with gross return $R>0$; there is no storage,
productive investment, government debt or outside asset. The young may
promise at most $D$ goods next period. Assume
\[
 0<D<\frac{Y}{1+\beta+\beta^2}.                                   \tag{1}
\]
In a stationary equilibrium their borrowing is $d/R$ with repayment
$d\leq D$. Middle-aged saving is $s\geq0$. Lifetime budgets are
$c_y=d/R$, $c_m=Y-d-s$, and $c_o=Rs$. Loans are enforceable. Goods and
loan markets clear and all cohort plans are optimal at the same rate.
These are full-employment endowments; no inference about unemployment is
built into the asset-market equilibrium.

\begin{theorem}[Unique stationary natural rate and resource clearing]
Under (1), the borrowing limit binds, middle-aged saving is
\[
 s=\frac{\beta}{1+\beta}(Y-D),
\]
and the unique stationary gross natural rate is
\[
 R^*=\frac{gD(1+\beta)}{\beta(Y-D)}.                              \tag{2}
\]
At this rate, consumption is positive and total consumption equals the
middle-aged generation's endowment. The log natural rate $r^*=\log R^*$
may be negative even though the economy has no missing goods.
\end{theorem}
\begin{proof}
Given repayment $d\in(0,Y)$, strict concavity of
$\log(Y-d-s)+\beta\log(Rs)$ gives the unique saving choice
$s=\beta(Y-d)/(1+\beta)$. Substitution into lifetime utility leaves,
up to constants independent of $d$,
$\log d+\beta(1+\beta)\log(Y-d)$. Its strictly decreasing derivative
crosses zero at $d=Y/(1+\beta+\beta^2)$. Assumption (1) therefore makes
$d=D$ the unique constrained optimum for the young. At date $t$, let the
middle cohort have size $N$; there are $gN$ young and $N/g$ old.
Loan clearing requires $Ns=gND/R$, which uniquely gives (2).

The young consume $D/R$, the middle aged $Y-D-s$, and the old receive
$Rs$ on the previous period's stationary savings. Total consumption per
middle-aged household is
\[
 gD/R+Y-D-s+(R/g)s=Y,
\]
because loan clearing gives $gD/R=s$ and $Rs/g=D$.
All terms are strictly positive by (1) and $R^*>0$. Equation (2) can be
below one, for instance by choosing sufficiently small positive $D$;
this does not change the resource identity.
\end{proof}
A small borrowing limit has two reinforcing effects: less young demand
for loans and, in the new stationary state, more disposable middle-age
income available for saving. Both are part of the source mechanism.
Within this endowment specialization, ``investment'' is zero and net
desired saving is saving by lenders minus dissaving by borrowers.

\section{Comparative statics, bounds and a lower-bound test}
At an arbitrary gross rate define excess net desired saving per
middle-aged household as
\[
 E(R)=\frac{\beta}{1+\beta}(Y-D)-\frac{gD}{R}.                     \tag{3}
\]
It is strictly increasing, with its sole root (2). Thus $E(R_0)>0$ if
and only if $R^*<R_0$.

\begin{proposition}[Identified signs and a transparent identified interval]
Within the maintained model,
\begin{align*}
 d\log R^*&=d\log g+
 \left(\frac1D+\frac1{Y-D}\right)dD
 -\frac{dY}{Y-D}-\frac{d\beta}{\beta(1+\beta)}.                    \tag{4}
\end{align*}
Suppose the admissible primitive set is the Cartesian box
$g\in[g_-,g_+]$, $D\in[D_-,D_+]$, $Y\in[Y_-,Y_+]$ and
$\beta\in[\beta_-,\beta_+]$, with positive endpoints, $\beta_+<1$,
and every point satisfying (1). The sharp set of gross rates is the
interval $[R_-,R_+]$, where
\[
 R_- =\frac{g_-D_-(1+\beta_+)}{\beta_+(Y_+-D_-)},\qquad
 R_+ =\frac{g_+D_+(1+\beta_-)}{\beta_-(Y_--D_+)}.                  \tag{5}
\]
With fixed expected log inflation $\bar\pi$, all admissible rates violate
the zero-nominal-rate bound exactly when $\log R_+<-\bar\pi$; none
violates it exactly when $\log R_-\geq-\bar\pi$.
\end{proposition}
\begin{proof}
Taking logarithms of (2) and differentiating gives (4). Thus $R^*$ is
increasing in $g,D$ and decreasing in $Y,\beta$. The extrema on the box
are the displayed corners and are attainable. The box is connected and
$R^*$ continuous, so its image in $\R$ contains every intermediate
value, proving sharpness in this explicitly unrestricted box class.
The two lower-bound claims follow by comparing the interval endpoints
with $\exp(-\bar\pi)$ and remembering the strict inequality defining a
violation.
\end{proof}
If data restrict primitives jointly, a box can be only an outer region,
and its corner bounds cease to be sharp. An observed yield alone does not
identify the mechanism: for fixed $Y,\beta$ and any observed $R>0$,
\[
 g=\frac{R\beta(Y-D)}{D(1+\beta)}                                 \tag{6}
\]
gives the same rate for every $D$ satisfying (1). These are different
valid equilibria, not relabelings. Age-specific debt and consumption data
can distinguish them; rate data cannot.

A full-employment rate below the monetary bound establishes the
infeasibility of implementing \emph{that} real allocation through the
specified nominal rate and inflation expectations. Persistent output
shortfalls additionally require nominal rigidities, policy and expectation
rules; they do not follow from (2) alone. Safe-asset premia likewise require
more than the single enforceable asset used here.

\section{Nonlinear attribution depends on the path}
Fix $g,\beta$ and vary $(Y,D)$ inside the region (1). Along any continuously
differentiable admissible path, the income and debt contributions in (4)
sum to the total log-rate change. Consider $Y_1>Y_0$ and $D_1>D_0$
with the whole rectangle admissible.

\begin{proposition}[Exact order sensitivity]
The debt contribution when debt changes before income exceeds the debt
contribution when debt changes after income by
\[
 \log\!\left(
 \frac{(Y_0-D_0)(Y_1-D_1)}{(Y_0-D_1)(Y_1-D_0)}\right)>0.           \tag{7}
\]
The income contribution changes by the opposite amount, leaving the
sum invariant.
\end{proposition}
\begin{proof}
At fixed $Y$, the debt-induced log-rate change is
$\log(D_1/D_0)+\log[(Y-D_0)/(Y-D_1)]$. Evaluate it at $Y_0$ and
$Y_1$ and subtract to obtain (7). Positivity also follows by integrating
$1/(Y-D)$ over $D\in[D_0,D_1]$: this integrand is strictly larger at
$Y_0$ than at $Y_1$. The total rate change is the difference of a
single-valued function at the endpoints, so the residual income
contributions differ by the negative of (7).
\end{proof}
Thus a share assigned to deleveraging is conditional on the counterfactual
path even when every primitive is known. Averaging orders is a possible
reporting convention, but it cannot identify an unobserved primitive
change or remove the need to declare the convention.

\section{World aggregation and external balances}
For a stationary integrated world with countries $j$, impose a common
cohort growth factor $g$ so that country population shares are constant.
Country $j$ has middle-aged population $N_j$, primitives $Y_j,D_j,\beta_j$
satisfying (1), and enforceable international loans. Saving is
$s_j=\beta_j(Y_j-D_j)/(1+\beta_j)$. A single world rate clears at
\[
 R_W=\frac{g\sum_jN_jD_j}{\sum_jN_js_j}.                           \tag{8}
\]
Country end-of-period net positions in newly issued one-period claims are
$F_j=N_j(s_j-gD_j/R_W)$, with $\sum_jF_j=0$.

\begin{proposition}[A national claim position is not a world excess]
Equation (8) is the unique stationary world asset-market equilibrium.
A country can have positive $F_j$ at that equilibrium, but a positive
world sum of these same positions is impossible. If the primitive vectors
are fixed, interpreting one country's $s_j-gD_j/R_W$ as global excess
saving produces an incorrect natural-rate equation.
\end{proposition}
\begin{proof}
Each household's optimality argument is unchanged under international
trade in identical claims. Summing desired loan supply and demand gives
$\sum N_js_j=g\sum N_jD_j/R$, whose unique positive solution is
(8). Substitution proves $\sum F_j=0$; individual summands can have
opposite signs when saving-to-borrowing ratios differ. Summing all household
budgets cancels both domestic and international claims, so world goods
clear. A single country's imbalance is therefore an external position,
not an additional world resource or an equilibrium violation.
These new-claim positions are not generally current accounts or changes
in net foreign assets: maturing foreign claims must also be included
when calculating those flows.
\end{proof}
The model leaves investment, asset safety, bequests, mortality, productivity
and distributional heterogeneity largely aside. Its sharp interval and
order-sensitivity results are restricted answers, not causal attribution
of observed global yields. The empirical next step is a joint primitive
region from age-specific balance sheets and cross-country positions,
followed by robustness to richer asset supply and nominal regimes.

\section{Completion Estimate}
40\%

This is a post hoc, heuristic estimate for the full catalogue target.
The restricted overlapping generations economy has a unique natural rate, sharp primitive box bounds, exact path sensitivity, and correct world asset clearing. Causal attribution of observed rate declines, productive investment, safe asset premia, heterogeneous saving, and monetary mechanisms sustaining output shortfalls remain outside this endowment benchmark.

\begin{thebibliography}{9}
\bibitem{em} G. B. Eggertsson and N. R. Mehrotra (2014).
\emph{A Model of Secular Stagnation}. NBER Working Paper 20574,
October 2014. Author-hosted 53-page original consulted, Section 2,
budget constraints (1)--(4), asset-market equilibrium and Section 2.1.
This is the two-author working paper, not the later quantitative
three-author journal article.
\url{https://gautieggertsson.org/research/a-model-of-secular-stagnation/w20574.pdf}.
\end{thebibliography}

\end{document}
